\begin{textsample}{POS Dim 3 – persona\_gemini – Score 17.00 – t596\_gemini.txt}  \label{ex:f3_pos_047}
To address the climate crisis, we must phase out diesel and petrol vehicles within the next decade.
In Europe, the \textbf{transport} \textbf{sector} is responsible for over a quarter of all CO2 \textbf{emissions}, and road \textbf{transport} alone accounts for 45 % of that.
Research conducted by the Ecologic Institute and Greenpeace Germany has identified six key policy \textbf{measures} to accelerate this transition.
The first is to implement a ban on the \textbf{sale} of new internal combustion engine \textbf{cars} while mandating \textbf{production} quotas for electric vehicles ( EVs ).
This must be supported by the expansion of a subsidized charging \textbf{infrastructure} to make EVs a viable \textbf{option} for everyone.
Furthermore, policies should be flexible, with incentives that adapt to changing \textbf{market} \textbf{trends}.
A mix of policies, including both subsidies and mandates, is necessary for a successful transition.
This approach should involve \textbf{taxing} polluters while offering incentives for the adoption of EVs.
Crucially, the benefits of this \textbf{shift} must be clearly communicated to the \textbf{public}.
It is important to recognise that electric vehicles are part of a much broader \textbf{solution}.
The ultimate goal is a systemic \textbf{shift} towards reducing the overall \textbf{number} of personal \textbf{cars} on our roads.
This involves boosting \textbf{public} \textbf{transport} powered by renewable \textbf{energy}, and making it easier and safer for people to choose \textbf{cycling} and walking.

% matched lemmas: car, cycle, emission, energy, infrastructure, market, measure, number, option, production, public, sale, sector, shift, solution, tax, transport, trend
\end{textsample}
